平方函数论证需要一个一致的 $L^p$ 界，来控制作为伸缩算子输入的一列光滑环带截断函数。

\begin{Proposition}[环带 Littlewood--Paley 界]\label{ext:stein-annular-square}
设 $n\ge1$、$1<p<\infty$，令 $q\in C_c^\infty(\mathbb R^n)$ 支撑于 $\{c\le|\xi|\le C\}$，其中 $0<c\le C<\infty$。采用
\[
 \widehat f(\xi)=\int_{\mathbb R^n}e^{-ix\cdot\xi}f(x)\,dx,\qquad
 \check q(x)=(2\pi)^{-n}\int_{\mathbb R^n}e^{ix\cdot\xi}q(\xi)\,d\xi .
\]
对 $j\in\mathbb Z$，在 Schwartz 函数上定义 $P_j=q(2^{-j}D)$，公式为 $\widehat{P_j f}(\xi)=q(2^{-j}\xi)\widehat f(\xi)$。在 $L^p(\mathbb R^n)$ 上使用卷积公式 $P_j f=2^{jn}\check q(2^j\cdot)*f$。则
\[
 \left\|\left(\sum_{k\in F}|P_{k+\ell} f|^2\right)^{1/2}\right\|_{L^p(\mathbb R^n)}
 \le C_{n,p,q}\|f\|_{L^p(\mathbb R^n)}
\]
对每个有限 $F\subset\mathbb Z$、每个整数移位 $\ell$ 和每个 $f\in L^p(\mathbb R^n)$ 成立。该界也适用于全和。未经伸缩的核 $\check q$ 为 Schwartz 函数且积分为零，因为 $q(0)=0$。它及其一阶导数的衰减给出 \cite{r-32}, Theorem~6.1.2, (6.1.3) 中的核界，因此 (6.1.4), p.~421 适用。限制到指标 $F+\ell$ 仍保持该界。
\end{Proposition}
